\documentclass[aspectratio=169]{beamer}

\usepackage{fontspec}
\usepackage{polyglossia}
\usepackage{color}
\usepackage{listings}
\setdefaultlanguage{russian}
\usepackage{csquotes}
\usetheme{metropolis}
\metroset{background=light}
\metroset{sectionpage=none}
\metroset{numbering=fraction}

\definecolor{ssaublue}{RGB}{77,144,205}
\setbeamercolor{frametitle}{fg=white, bg=ssaublue}
\setbeamercolor{title separator}{fg=ssaublue}

\setmainfont[
    Path = ../fonts/elektratextpro/,
    Extension = .otf,
    UprightFont = elektratextpro,
    BoldFont = elektratextpro_bold,
    ItalicFont = elektratextpro_italic,
    BoldItalicFont = elektratextpro_bolditalic
]{elektratextpro}

\setsansfont[
    Path = ../fonts/elektratextpro/,
    Extension = .otf,
    UprightFont = elektratextpro,
    BoldFont = elektratextpro_bold,
    ItalicFont = elektratextpro_italic,
    BoldItalicFont = elektratextpro_bolditalic
]{elektratextpro}

\setmonofont[
    Path = ../fonts/elektratextpro/,
    Extension = .otf,
    UprightFont = elektratextpro,
    BoldFont = elektratextpro_bold,
    ItalicFont = elektratextpro_italic,
    BoldItalicFont = elektratextpro_bolditalic
]{elektratextpro}

\newfontfamily\cyrillicfont[
    Path = ../fonts/elektratextpro/,
    Extension = .otf,
    UprightFont = elektratextpro,
    BoldFont = elektratextpro_bold,
    ItalicFont = elektratextpro_italic,
    BoldItalicFont = elektratextpro_bolditalic
]{elektratextpro}

\newfontfamily\cyrillicfontsf[
    Path = ../fonts/elektratextpro/,
    Extension = .otf,
    UprightFont = elektratextpro,
    BoldFont = elektratextpro_bold,
    ItalicFont = elektratextpro_italic,
    BoldItalicFont = elektratextpro_bolditalic
]{elektratextpro}

\newfontfamily\cyrillicfonttt[
    Path = ../fonts/elektratextpro/,
    Extension = .otf,
    UprightFont = elektratextpro,
    BoldFont = elektratextpro_bold,
    ItalicFont = elektratextpro_italic,
    BoldItalicFont = elektratextpro_bolditalic
]{elektratextpro}

% Настройка листингов для Beamer
\lstset{
    basicstyle=\ttfamily\footnotesize,
    commentstyle=\color{gray},
    columns=fullflexible,
    keepspaces=true,
    breaklines=false,
    extendedchars=true,
    texcl=true,
    showstringspaces=false
}

\title{Веб-сервер, прокси и туннелирование}
\author{Правдин Иван}
\date{24.05.2025}

\begin{document}
    \maketitle

    \begin{frame}{Что такое веб-сервер?}
        \textbf{Веб-сервер} — программное обеспечение, которое принимает HTTP-запросы от клиентов, обрабатывает и возвращает HTTP-ответы
        
        \vspace{0.5cm}
        
        \textbf{Принцип работы:}
        \begin{enumerate}
            \item Клиент отправляет HTTP запрос
            \item Сервер обрабатывает запрос и возвращает HTTP ответ
        \end{enumerate}
        
        \vspace{0.5cm}
        
        \textbf{Примеры веб-серверов:}
        \begin{itemize}
            \item \textbf{Apache}
            \item \textbf{Nginx}
            \item \textbf{Python http.server}
        \end{itemize}
    \end{frame}

    \begin{frame}{Демонстрация: запустим простой веб-сервер}
        \textbf{Что покажем:}
        \begin{itemize}
            \item Запуск статического и динамического веб-сервера
            \item Обращение к веб-серверу из браузера, developers tools
            \item Обращение к веб-серверу с помощью curl
        \end{itemize}
    \end{frame}

    \begin{frame}{Прокси}
        \textbf{Прокси сервер} — сервер-посредник между клиентом и сервером
        
        \vspace{0.5cm}
        
        \begin{columns}[t]
            \begin{column}[t]{0.48\textwidth}
                \textbf{Forward proxy:}
                \begin{itemize}
                    \item Клиент → Прокси → Интернет
                    \item Скрывает клиента от сервера
                    \item Пример: корпоративный прокси
                \end{itemize}
            \end{column}
            
            \begin{column}[t]{0.48\textwidth}
                \textbf{Reverse proxy:}
                \begin{itemize}
                    \item Интернет → Прокси → Сервер
                    \item Скрывает сервер от клиента
                    \item Пример: nginx перед приложением
                \end{itemize}
            \end{column}
        \end{columns}
    \end{frame}

    \begin{frame}{Демонстрация: nginx в качестве reverse proxy}
        \textbf{Что покажем:}
        \begin{itemize}
            \item Маршрутизация по портам
            \item Маршрутизация по path
            \item Маршрутизация по заголовку Host
        \end{itemize}
    \end{frame}

    \begin{frame}{Туннелирование}
        Туннелирование — технология инкапсуляции одного сетевого протокола в другой для создания канала передачи данных
        
        \vspace{0.5cm}
        
        \textbf{Назначение:}
        \begin{itemize}
            \item Обход сетевых ограничений (NAT, firewall)
            \item Обеспечение безопасности передачи данных
            \item Соединение изолированных сетей
        \end{itemize}
        
        \vspace{0.5cm}
        
        \textbf{Принцип работы:}
        \begin{enumerate}
            \item Исходный трафик упаковывается в другой протокол
            \item Передается через промежуточный сервер
            \item Распаковывается на стороне получателя
        \end{enumerate}
    \end{frame}

    \begin{frame}{Демонстрация: Туннелирование через SSH}
        \textbf{Что покажем:}
        \begin{itemize}
            \item Запуск веб-сервера на удаленной машине
            \item Создание SSH туннеля: локальный порт → удаленный сервис
            \item Доступ к удаленному сайту через локальный адрес
        \end{itemize}
    \end{frame}

    \begin{frame}{FRP: reverse proxy + тунель}
        \textbf{Fast Reverse Proxy (FRP)} — обратный прокси сервер с поддержкой туннелирования до клиентов
        
        \vspace{0.5cm}
        
        \textbf{Возможности FRP:}
        \begin{itemize}
            \item HTTP туннели
            \item TCP туннели
            \item UDP туннели
            \item Работа через NAT и firewall
        \end{itemize}
        
        \vspace{0.5cm}
        
        \textbf{Архитектура:}
        \begin{center}
            Клиент → FRP Server ← FRP Client ← Сервер
        \end{center}
        
    \end{frame}

    \begin{frame}{Демонстрация: FRP в действии}
        \textbf{Что покажем:}
        \begin{itemize}
            \item Установка и настройка FRP клиента
            \item Создание HTTP туннеля для веб-сервера
            \item Создание TCP туннеля для SSH доступа
            \item Доступ к внутренним сервисам через внешний домен
        \end{itemize}
    \end{frame}

    \begin{frame}{Что запомнить}
        \begin{itemize}
            \item \textbf{HTTP} — основа веб-коммуникации (заголовки, статус коды)
            \item \textbf{Forward Proxy} — клиент → прокси → интернет
            \item \textbf{Reverse Proxy} — интернет → прокси → сервер
            \item \textbf{Туннелирование} — создание канала передачи данных через сеть
        \end{itemize}

        \vspace{1cm}

        \begin{center}
            \textbf{Вопросы?}
        \end{center}
    \end{frame}

\end{document}